\documentclass[11pt]{article}
\usepackage[a4paper,margin=1in]{geometry}
\usepackage{amsmath,amssymb,mathtools,bm}
\usepackage{enumitem,booktabs,array}
\usepackage{tcolorbox}
\usepackage{hyperref}
\usepackage[protrusion=true,expansion=false]{microtype}
\hypersetup{colorlinks=true,linkcolor=blue,urlcolor=blue,citecolor=blue}
\setlist[itemize]{topsep=2pt,itemsep=2pt}
\setlist[enumerate]{topsep=2pt,itemsep=2pt}
\newtcolorbox{keybox}{colback=blue!3!white,colframe=blue!40!black,boxrule=0.4pt,arc=1mm}
\newtcolorbox{alertbox}{colback=red!3!white,colframe=red!40!black,boxrule=0.4pt,arc=1mm}
\newtcolorbox{answerbox}{colback=green!3!white,colframe=green!40!black,boxrule=0.4pt,arc=1mm}
\newtcolorbox{drillbox}{colback=yellow!5!white,colframe=orange!60!black,boxrule=0.4pt,arc=1mm}
\newcommand{\EV}{\mathbb{E}}
\newcommand{\Var}{\mathrm{Var}}
\newcommand{\Cov}{\mathrm{Cov}}
\newcommand{\blank}[1]{\underline{\hspace{#1}}}

\title{W13-03: Krueger, Sautner \& Starks (2020, RFS) \\
\large ``The Importance of Climate Risks for Institutional Investors'' --- Active Recall Workbook}
\author{Banking \& Risk Management --- Exam Prep}
\date{\today}
\begin{document}\maketitle

\begin{keybox}
\textbf{Paper in one sentence.} In a global survey of $\sim$440 institutional investors (2017--2018), KSS document that investors view climate risks as \emph{financially material} --- especially regulatory and physical risks --- and cite \emph{risk management}, not ethics, as the primary motive for engagement; divestment is a minority strategy.

\textbf{Placement.} The canonical survey of institutional-investor climate attitudes. Companion to Dyck-Lins-Roth-Wagner (W13-04) on whether institutions actually move firm CSR in the cross-section.
\end{keybox}

\section*{Round 1: Complete Walk-Through}

\subsection*{1.1 Survey design}
Global survey: $\sim$440 institutional investors (mutual funds, pensions, endowments, sovereign wealth funds). Fielded 2017--2018. Questions cover materiality, motives, strategies (engagement vs.\ divestment), and time horizons.

\subsection*{1.2 Headline findings}
\begin{itemize}
\item $\sim$60--70\% of institutions say climate risks are financially material.
\item Regulatory risk ranks highest, followed by physical risk, then technological.
\item Main motive: \emph{risk management}; only a minority cite ethics.
\item Preferred strategy: \emph{engagement} (dialogue, proxy voting) over divestment.
\item Horizons: many say $>5$ years; some $>10$ years.
\item Larger institutions, European institutions, and ESG-mandate funds are most climate-aware.
\end{itemize}

\subsection*{1.3 Cross-sectional variation}
\begin{itemize}
\item EU $>$ US, Asia-Pacific.
\item Large $>$ small; signatories to the Principles for Responsible Investment are much more active.
\item Pensions $>$ hedge funds.
\end{itemize}

\subsection*{1.4 Policy implications}
\begin{itemize}
\item Institutional pressure is a real force in corporate climate decisions.
\item Disclosure-regulation (TCFD, SFDR, EU Taxonomy) aligns with stated institutional demands.
\item Institutions seek \emph{risk-adjusted} returns; climate is priced as risk, not moral consumption.
\end{itemize}

\subsection*{1.5 Caveats}
\begin{alertbox}
\begin{itemize}
\item Self-report: strategic answers possible (esp.\ on ESG).
\item Non-response bias.
\item Survey cross-section; cannot test causality on portfolio changes.
\end{itemize}
\end{alertbox}

\section*{Round 2: Level 1 Fill-in}
\begin{enumerate}
\item Survey size: $\sim$\blank{1cm} institutions; years \blank{1cm}--\blank{1cm}.
\item Most-cited risk type: \blank{2cm}.
\item Primary motive: \blank{2cm}.
\item Preferred strategy: \blank{2cm}.
\item Geography: \blank{1cm} most climate-aware.
\end{enumerate}

\section*{Round 3: Level 2 Prompts}
\begin{enumerate}
\item Discuss regulatory, physical, and technological climate risks from an institutional perspective.
\item Why is engagement preferred over divestment?
\item Relate KSS to BBGL (2022) on belief heterogeneity.
\item Design a follow-up test on portfolio outcomes consistent with KSS.
\end{enumerate}

\section*{Round 4: Minimal Scaffolding}
From a blank page: (i) survey, (ii) materiality ranking, (iii) motives, (iv) strategies, (v) caveats.

\section*{Round 5: Blank Page}
Essay: climate risk through the institutional lens --- evidence from KSS.

\vspace{6pt}
\begin{drillbox}
\textbf{Speed Drill.} (1) Survey size. (2) Top risk type. (3) Main motive. (4) Preferred strategy. (5) Top regional cohort.
\end{drillbox}

\section*{Quick Reference \& Answer Key}
\begin{answerbox}
\begin{itemize}
\item \textbf{Survey.} $\sim$440 institutions, 2017--18.
\item \textbf{Materiality.} Regulatory $>$ physical $>$ technological.
\item \textbf{Motive.} Risk management; ethics secondary.
\item \textbf{Strategy.} Engagement preferred over divestment.
\end{itemize}
\end{answerbox}

\begin{keybox}
\textbf{30-second exam pitch.} Krueger, Sautner \& Starks (2020) survey about 440 global institutional investors in 2017--18 on climate risk. Most respondents (60--70\%) consider climate risks financially material, ranking regulatory risk highest, followed by physical and then technological risks. The dominant motive for action is \emph{risk management} rather than ethics; engagement with portfolio firms is strongly preferred to divestment; investment horizons of 5--10 years are common. European, large, and PRI-signatory investors are most climate-aware. The paper establishes that institutions treat climate risk as a first-order financial variable, aligning with disclosure regulations (TCFD, SFDR) and with Bernstein-Billings-Gustafson-Lewis (2022)'s belief-heterogeneity finding. It is the canonical survey reference for institutional climate attitudes.
\end{keybox}

\end{document}
